\pagebreak
\subsection{Conditional Register Load With Subtraction}
\verb|ldc DR,SR - IMM|
\subsubsection{Encoding}
This instruction will use the RRI format. 
The operation field (OP) wil be $0000_2$

\begin{center}
    \begin{tabular}{|B|B|B|B|B|B|B|B|B|B|B|B|B|B|B|B|B|B|B|}
        \hline
        \textbf{15} & \textbf{14} & \textbf{13} &
        \textbf{12} & \textbf{11} & \textbf{10} & \textbf{9} &
        \textbf{8} & \textbf{7} & \textbf{6} & \textbf{5} &
        \textbf{4} & \textbf{3} & \textbf{2} & \textbf{1} &
        \textbf{0} \\

        \hline
        \multicolumn{1}{|c|}{0} & 
        \multicolumn{1}{c|}{0} & 
        \multicolumn{1}{c|}{0} & 
        \multicolumn{1}{c|}{0} & 
        \multicolumn{2}{c|}{DST} &
        \multicolumn{2}{c|}{SRC} & 
        \multicolumn{8}{c|}{IMM} \\
        \hline
    \end{tabular}
\end{center}

\subsubsection{Operation}
$if (L = 1) \{ DR \leftarrow SR - SignExtend(IMM) \}$
\par 
If the link register is not set, skip this instruction.
Subtract the the sign-extended immediate (IMM) from the contents of the source register (SR).
Latch the overflow into the link register.
Place the result into the destination register (DR).

\subsubsection{Associated Pseudo-Instructions}
The assembler will provide `pseudo-instructions' for the programmer's convenience.

\begin{center}
    \begin{tabular}{|c|c|c|}
        \hline
        Name & Syntax & Equivalent \\
        \hline
        No-Operation & \verb|nop| & \verb|ldc zr, zr + 0|  \\
        \hline
        Add Link & \verb|adl| & \verb|ldc ac, ac - 1| \\
        \hline
        Jump Relative Conditional & \verb|jrc OFFSET| & \verb|ldc pc, pc - OFFSET| \\
        \hline
    \end{tabular}
\end{center}
